\documentclass{plantillaPPT}

\titulo{7}{Áreas Entre Regiones}

\begin{document}

\primeraDiapo

\begin{frame}{Cálculo de áreas entre regiones}
    \framesubtitle{Práctica 7}

1. Determinar el área de cada una de las siguientes regiones del plano. \\ \vspace{0.3cm}
\((a)\) La región bajo la parábola de la ecuación \(y=3-x^2\) y sobre la recta \(y=-1\). \\ \vspace{0.3cm}
\((b)\) El triángulo de vértices \((1,1), (5,2) \text{ y } (4,8)\). \\ \vspace{0.3cm}
\((c)\) La región encerrada entre las gráficas de \(x-y+1=0\) e \(y = x^3-3x^2-3x+1\)

\end{frame}

\begin{frame}{Cálculo de áreas entre regiones}
    \framesubtitle{Práctica 7}

2. Considere las siguientes regiones \(R_1\) y \(R_2\) del plano acotadas por las curvas.\vspace{-0.10cm}  \\  
\begin{center}
\(C_1: y=1 \qquad C_2:y=\sqrt{|x|} \qquad C_3: x^2+y^2=2\)
    \includegraphics[width=0.55\textwidth]{C2 2025-1/imagenes/ayud7img.png}
\end{center}
\vspace{-0.25cm}

\((a)\) Defina las regiones de forma conjuntista. \\
\((b)\) Calcule el área de la región \(R = R_1 \cup R_2\), integrando respecto a ambas variables.

\end{frame}

\begin{frame}{Problema de Restricción}
    \framesubtitle{Práctica 7}

3. Sean \(a \in \mathbb{R}\) y \(R\) la región definida por:
\[R = \{(x,y)\in \mathbb{R}^2 : 1\leq x+y\leq a,\; x,y\geq0\}\]
Determine, si existe, el valor de \(a\) de modo que la medida del área de la región sea \(12u^2\)

\end{frame}

\end{document}